
\section{\texorpdfstring{Description of the participating subjects
}{Description of the participating subjects }}\label{description-of-the-participating-subjects}

\subsection{\texorpdfstring{\ul{Principal
Investigator}}{Principal Investigator}}\label{principal-investigator}

\subsubsection{Scientific Contributions}\label{scientific-contributions}

I have been in the research field for more than 15 years, and my
scientific publications are enclosed in my CV. I am an Associate Editor
for RCIM the most prestigious journal paper in the field of applied
robotics science, edited from Elsevier. I have about ten patents
deposited, and one international patent has been assigned to a company
and will be active until 2037.

I coordinated or was responsible for 17 projects, of which 8 were EU
projects (FP7, H2020, HE), 2 EU meta-projects, one UK collaboration, one
national, two regionals, and 3 Research and Innovation Contracts with
Companies. The 18 projects I managed had a financial capacity of just
over €6M (max €810k, min €15k). I participated as a unit member in 12
projects (WP Leader and/or Task Leader).

My scientific production has focused on developing methodologies to
extend the ``capabilities'' of industrial robots so they can be used in
unstructured environments, significantly increasing the added value of
the production system and improving the quality of work for human
operators. My contributions have focused on three enabling technologies
for industrial robots (only the most essential products are listed):
motion and task planning, interaction control, modeling, and robotized
machining.

In addition to contributions in these areas, I investigated the use of
anthropomorphic robots for rehabilitation and the development of new
mechatronic devices and new robots for cranial or spinal surgery
{[}AJ44{]}. Another area where I sought to make original contributions
was the development of systems to make industrial robotic applications
safe, even without barriers. Specifically, we investigated how certified
and non-certified sensor networks could be used. This activity is
characterized by a multidisciplinary solid involving the CNR-IEIIT
institute.

Some works have been more related to new frameworks, and only recently
have I started investigations into the use of Deep and Reinforcement
Learning in the industrial sector. The achievement of the scientific
results described is linked to both scientific and managerial
methodological skills.

Regarding methodology, the results demonstrate that I possess (i)
scientific solid rigor, (ii) analytical and synthetic skills in
multi-body system modeling, (iii) a propensity for multidisciplinary,
(iv) rigorous study of state-of-the-art

Regarding managerial skills, the results show that I possess (i)
excellent project management skills, (ii) strong networking and lobbying
capabilities, (iii) administrative and managerial skills acquired
through the management of European projects, (iv) attention to the
development of the human resources I manage. These managerial skills
have been fundamental for me in managing three laboratories in parallel
and four PhDs. The two aspects on which I have focused my attention have
been (1) very detailed activity planning and (2) formalization of
knowledge within the labs, explicitly defining methods for developing
software and mechatronic projects that enable sharing and maintenance.

\subsubsection{Technological Transfer}\label{technological-transfer}

The transfer and application of my research results for the benefit of
the scientific community over the years have followed three operational
models:

\begin{enumerate}
\def\labelenumi{(\Alph{enumi})}
\item
  participation in European projects and exploitation of research
  results.
\item
  transfer of intellectual property to industrial operators.
\item
  sharing of developed code in applications through Free and Open-Source
  Software (FOSS) principles.
\end{enumerate}

Regarding point (A), the three EU projects for which I was the Principal
Investigator for CNR have received awards for achievements in technology
transfer. Specifically, in ShareWork, the European
Commission\textquotesingle s Innovation Radar selected three
implementation results developed by CNR
\href{https://innovation-radar.ec.europa.eu/resultbykeyword/sharework}{(https://innovation-radar.ec.europa.eu/resultbykeyword/sharework)}.
This Award recognizes the impact of CNR\textquotesingle s scientific
results within the project. As for the Robofoot project {[}PC14{]} was
ranked among the top five projects of the Seventh Framework Programme
(FP7-EU) for impact and implementation quality. To maximize the
exploitation of results, my goal in recent years has been to seek
funding for high-TRL actions, such as Large-Scale Demonstrators or
actions related to translational research activities (e.g., leveraging
Horizon Joint Undertaking calls). This choice is motivated by both
cultural and managerial reasons. The Italian industrial fabric requires
research entities to serve as benchmarks for the growth of the
manufacturing sector. Guiding these companies through transformations of
production systems by demonstrating high-impact project results creates
added value for the country. From a managerial perspective, the high
preparation of my group\textquotesingle s human resources and the
availability of robotic solutions accumulated over a decade of European
projects allows us to realize industrial demonstrators relatively
quickly, reducing startup times and energy for new initiatives. This, in
turn, provides a strong financial lever for investing in PhDs and basic
research.

I support knowledge sharing through Free and Open-Source Software (FOSS)
paradigms (C). This approach allows the creation of large, collaborative
working communities. My research group has activated 70 public
Open-Source projects, while the CARI lab
(\href{https://github.com/JRL-CARI-CNR-UNIBS}{https://github.com/JRL-CARI-CNR-UNIBS)}
has 84 active projects with an average of 50 users. The FOSS model
offers several advantages, such as being more transparent,
cost-effective, and robust and enabling rapid improvements through
community collaboration. In this field, the open-source community formed
around ROS (Robotic Operative System) has been a catalyst for
transforming research worldwide.

\subsubsection{Industrial and Institutional
Roles}\label{industrial-and-institutional-roles}

Regarding products and scientific qualifications, the most relevant
aspects are consulting technical-scientific collaboration and support
for companies. Specifically, I initiated industrial contracts worth
approximately €600k. A long collaboration track with COMAU has been set
in the last ten years.

I have participated in 12 evaluation commissions as an expert. Regarding
my participation in evaluation committees, I have served on the
monitoring committee of Lombardy Region overseeing the EcoCirc project;
I am a CNR expert for the evaluation of MISE projects, and I have been
called upon as a member to evaluate EU-funded projects as well as
international projects. I have served as the final evaluator for
European initiatives.

At the national level, I have also participated in organizing some
courses for SIRI (Italian Society of Industrial Robotics,
\href{http://www.siri.it}{www.siri.it}). I was invited to speak at the
MECPSE fair, a national reference for automation. Additionally, for
MECPSE, I was a member of the committee that awarded the Innovation
Awards in 2021. At the international level, I was invited to participate
in several workshops organized by the European Robotic Forum,
particularly in robotics-related events for advanced manufacturing. I
was also invited to speak at the first B2B European robotics event.

In 2005, I founded PoliBrixia (\href{https://www.polibrixia.it}{www.polibrixia.it}), a spin-off of the
University of Brescia. This company is still active today and is a
dynamic player in Research and Innovation in Industry 4.0 and
Rehabilitation Robotics. I sold my shares when I became a permanent
researcher at CNR-ITIA in 2011.

\subsection{\texorpdfstring{Host Institution
}{Host Institution }}\label{host-institution-4}

COMAU (COnsorzio MAcchine Utensili) is an Italian multinational company
specializing in industrial automation and robotics. Founded in 1973,
COMAU is headquartered in Turin, Italy. The company offers automation
solutions across various industries, such as automotive, aerospace,
renewable energy, heavy industry, and logistics.

\textbf{Key Areas of Expertise}:

\begin{itemize}
\item
  Robotics: COMAU is a global leader in designing and producing
  industrial robots for welding, assembly, material handling, and
  painting. Their robots are used in high-precision applications and can
  be customized for different industries.
\item
  Industrial Automation: The company provides automation solutions,
  including assembly lines, automated manufacturing systems, and
  equipment for body-in-white (the stage of car manufacturing where the
  car body is assembled).
\item
  Digital Manufacturing: COMAU integrates digital technologies such as
  IoT, AI, and data analytics into manufacturing processes, supporting
  Industry 4.0 initiatives to optimize production, maintenance, and
  resource management.
\item
  Sustainability: COMAU creates solutions contributing to
  sustainability, including automation systems for renewable energy
  projects, particularly in wind energy, solar energy, and energy
  storage.
\end{itemize}

COMAU operates in numerous countries worldwide, including major markets
in Europe, North and South America, and Asia. It has over 4,000
employees and a network of global innovation centers, offices, and
production plants.

COMAU is known for its cutting-edge technology, innovation in robotics,
and contributions to modernizing manufacturing and production processes
across industries.

The track of innovation of COMAU is also demonstrated by the large
number of EU projects in which it has been involved. As shown at
\url{https://dashboard.tech.ec.europa.eu/}, COMAU has been involved in 46 EU
projects from FP7, ranking the company at 111 positions over 5583 funded
institutions in Europe from FP7. 34\% of the project is ongoing, making
COMAU one of Europe\textquotesingle s most active Robotics companies in
research.

COMAU gathered about 17M€ from FP7 from participating in such projects,
which has allowed COMAU to be among the most advanced worldwide robotic
innovators.

COMAU has a well-consolidated knowledge of the industrial robotics field
and will contribute to the project with its products, laboratories, and
research facilities, where advanced solutions are developed and
extensively tested. Moreover, COMAU will provide support and competence
in integrating new products and processes by adopting the latest
hardware and software technologies.

\subsection{\texorpdfstring{Affiliation Institution and Research
Organization
}{Affiliation Institution and Research Organization }}\label{affiliation-institution-and-research-organization}

The National Research Council of Italy is the largest national public
research organization, nation-wide distributed in a network of 80+
Institutes, cross-disciplinary fostering science and innovation in
biotechnology, medicine, materials, environment and land, information
and communications, advanced systems of production, judicial and
socio-economic sciences, classical studies and arts.

CNR 8000+ people are committed to promoting, transferring, and improving
research activities in the main sectors of knowledge and applications
for the scientific, technological, economic, and social development of
the Country.

Two different Institutes of CNR will follow the project:

\begin{itemize}
\item
  \textbf{STIIMA} - Institute of Intelligent Industrial Systems and
  Technologies for Advanced Manufacturing. Specifically, STIIMA conducts
  scientific research, development, technology transfer, training, and
  strategic roadmapping activities to contribute to
  enterprises\textquotesingle{} innovation, competitiveness, and
  sustainability and promote people\textquotesingle s central role in
  enterprises and society. In particular, the activities aim to design
  intelligent systems, enabling industrial technologies, products, and
  processes that co-evolve dynamically over time to meet different
  social and market needs and support new production paradigms.
  Alongside the Institute\textquotesingle s manufacturing expertise, it
  is clear that areas of primary interest are represented by the study
  of natural and synthetic fibers and their characterization,
  functionalization, and transformation processes, the development of
  solutions for ambient intelligence and ambient assisted living, the
  central role of the individual in the society of the future, with
  emphasis on the issue of assistance and the social and labor inclusion
  of vulnerable groups, advanced applications in the field of
  agricultural robotics and non-destructive quality control in the
  agri-food sector, bioinformatics and ecological informatics and the
  circular economy issue. CNR-STIIMA is the reference CNR institute for
  production systems (factory automation and robotics, mechatronics,
  industrial and service robotics, microsystems, controls, diagnosis,
  monitoring and new generation supervising systems, virtual design
  methods of products, virtual and augmented reality, process and
  factory design, management and business paradigms). CNR-STIIMA is a
  major partner of the European Technological Platform ManuFuture
  initiative (STIIMA\textquotesingle s Director is a member of the
  High-Level Group and of the Implementation Support Group), finalized
  in developing added high-value manufacturing based on research and
  innovation. CNR-STIIMA is a founding member of the EFFRA association
  (private party of Factories of the Future PPP), the Director being a
  member of the Industrial Research Advisory Group and the Ad-hoc
  Industrial Advisory Group. CNR-STIIMA is a founding member of
  euRobotics AISBL (a private party of SPARC).
\item
  \textbf{ISTC} Institute for Cognitive Science and Technology. The
  Institute participates in the project through the Planning and
  Scheduling Technology Lab (PSTLab - \url{http://istc.cnr.it/group/pst}), a
  group founded in 1997 and focused on artificial intelligence
  techniques for automated and interactive problem solving, integrated
  planning and scheduling algorithms, effective scheduling heuristics,
  dynamic scheduling, constraint programming, user interactive
  mixed-initiative systems. The group has been involved for over ten
  years in direct collaboration with the European Space Agency (ESA) to
  deliver several mission operational software systems (MEXAR, RAXEM,
  MrSPOCK) and two software infrastructures that are currently Agency
  assets (APSI, a planning software platform, and GOAC, a robot control
  software for autonomy). Between 2002 and 2006, the group coordinating
  entity of the RoboCare project, an initiative of the Italian Ministry
  of Research and University, was one of the first projects in Europe on
  robotic companions for older adult assistance. Since the start of FP7,
  the PSTLab has broadened its research interests toward new application
  areas, acquiring competencies in robotics (particularly human-robot
  interaction), user evaluation (for example, older adult users),
  ambient assisted living (smart environments integrating sensors,
  middleware, and intelligent services) and cognitive systems for
  training. The group has also worked in the manufacturing domain on
  national and international projects.\textbf{\hfill\break}
\end{itemize}